\section{Closed immersions of ringed spaces}
\label{section-closed-immersion}

\noindent
When do we declare a morphism of ringed spaces
$i : (Z, \mathcal{O}_Z) \to (X, \mathcal{O}_X)$
to be a closed immersion?

\medskip\noindent
Motivated by the example of a closed immersion of normal topological spaces
(ringed with the sheaf of continuous functors), or differential manifolds
(ringed with the sheaf of differentiable functions), it seems natural to
assume at least:
\begin{enumerate}
\item The map $i$ is a closed immersion of topological spaces.
\item The associated map $\mathcal{O}_X \to i_*\mathcal{O}_Z$
is surjective. Denote the kernel by $\mathcal{I}$.
\end{enumerate}
Already these conditions imply a number of pleasing results: For example
we prove that the category of $\mathcal{O}_Z$-modules is equivalent to
the category of $\mathcal{O}_X$-modules annihilated by $\mathcal{I}$
generalizing the result on abelian sheaves of
Section \ref{section-closed-immersions}.

\medskip\noindent
However, in the Stacks project we choose the
definition that guarantees that if $i$ is a closed immersion
and $(X, \mathcal{O}_X)$ is a scheme, then also $(Z, \mathcal{O}_Z)$
is a scheme. Moreover, in this situation we want $i_*$ and $i^*$
to provide an equivalence between the category of quasi-coherent
$\mathcal{O}_Z$-modules and the category of quasi-coherent
$\mathcal{O}_X$-modules annihilated by $\mathcal{I}$.
A minimal condition is that $i_*\mathcal{O}_Z$ is a
quasi-coherent sheaf of $\mathcal{O}_X$-modules.
A good way to guarantee that $i_*\mathcal{O}_Z$ is a
quasi-coherent $\mathcal{O}_X$-module is to assume that
$\mathcal{I}$ is locally generated by sections.
We can interpret this condition as saying ``$(Z, \mathcal{O}_Z)$ is
locally on $(X, \mathcal{O}_X)$ defined by setting some regular functions
$f_i$, i.e., local sections of $\mathcal{O}_X$, equal to zero''.
This leads to the following definition.

\begin{definition}
\label{definition-closed-immersion}
A {\it closed immersion of ringed spaces}\footnote{This is
nonstandard notation; see discussion above.} is a morphism
$i : (Z, \mathcal{O}_Z) \to (X, \mathcal{O}_X)$
with the following properties:
\begin{enumerate}
\item The map $i$ is a closed immersion of topological spaces.
\item The associated map $\mathcal{O}_X \to i_*\mathcal{O}_Z$
is surjective. Denote the kernel by $\mathcal{I}$.
\item The $\mathcal{O}_X$-module $\mathcal{I}$ is locally
generated by sections.
\end{enumerate}
\end{definition}

\noindent
Actually, this definition still does not guarantee that
$i_*$ of a quasi-coherent $\mathcal{O}_Z$-module is a
quasi-coherent $\mathcal{O}_X$-module. The problem is that
it is not clear how to convert a local presentation of
a quasi-coherent $\mathcal{O}_Z$-module into a local
presentation for the pushforward. However, the following
is trivial.

\begin{lemma}
\label{lemma-i-star-quasi-coherent}
Let $i : (Z, \mathcal{O}_Z) \to (X, \mathcal{O}_X)$
be a closed immersion of ringed spaces.
Let $\mathcal{F}$ be a quasi-coherent $\mathcal{O}_Z$-module.
Then $i_*\mathcal{F}$ is locally on $X$ the cokernel of
a map of quasi-coherent $\mathcal{O}_X$-modules.
\end{lemma}

\begin{proof}
This is true because $i_*\mathcal{O}_Z$ is quasi-coherent
by definition. And locally on $Z$ the sheaf $\mathcal{F}$
is a cokernel of a map between direct sums of copies
of $\mathcal{O}_Z$. Moreover, any direct sum of copies of the
{\it the same} quasi-coherent sheaf is quasi-coherent.
And finally, $i_*$ commutes with arbitrary colimits,
see Lemma \ref{lemma-i-star-right-adjoint}. Some details omitted.
\end{proof}

\begin{lemma}
\label{lemma-i-star-reflects-finite-type}
Let $i : (Z, \mathcal{O}_Z) \to (X, \mathcal{O}_X)$ be a morphism of ringed
spaces. Assume $i$ is a homeomorphism onto a closed subset of $X$ and
that $\mathcal{O}_X \to i_*\mathcal{O}_Z$ is surjective.
Let $\mathcal{F}$ be an $\mathcal{O}_Z$-module.
Then $i_*\mathcal{F}$ is of finite type if and only if
$\mathcal{F}$ is of finite type.
\end{lemma}

\begin{proof}
Suppose that $\mathcal{F}$ is of finite type.
Pick $x \in X$. If $x \not \in Z$, then $i_*\mathcal{F}$
is zero in a neighbourhood of $x$ and hence finitely generated in
a neighbourhood of $x$. If $x = i(z)$, then choose an open neighbourhood
$z \in V \subset Z$ and sections $s_1, \ldots, s_n \in \mathcal{F}(V)$
which generate $\mathcal{F}$ over $V$. Write $V = Z \cap U$ for some open
$U \subset X$. Note that $U$ is a neighbourhood of $x$. Clearly the
sections $s_i$ give sections $s_i$ of $i_*\mathcal{F}$ over $U$.
The resulting map
$$
\bigoplus\nolimits_{i = 1, \ldots, n} \mathcal{O}_U
\longrightarrow
i_*\mathcal{F}|_U
$$
is surjective by inspection of what it does on stalks
(here we use that $\mathcal{O}_X \to i_*\mathcal{O}_Z$
is surjective). Hence $i_*\mathcal{F}$
is of finite type.

\medskip\noindent
Conversely, suppose that $i_*\mathcal{F}$ is of finite type.
Choose $z \in Z$. Set $x = i(z)$. By assumption there exists
an open neighbourhood $U \subset X$ of $x$, and sections
$s_1, \ldots, s_n \in (i_*\mathcal{F})(U)$ which generate $i_*\mathcal{F}$
over $U$. Set $V = Z \cap U$. By definition of $i_*$ the sections
$s_i$ correspond to sections $s_i$ of $\mathcal{F}$ over $V$.
The resulting map
$$
\bigoplus\nolimits_{i = 1, \ldots, n} \mathcal{O}_V
\longrightarrow
\mathcal{F}|_V
$$
is surjective by inspection of what it does on stalks. Hence
$\mathcal{F}$ is of finite type.
\end{proof}

\begin{lemma}
\label{lemma-i-star-equivalence}
Let $i : (Z, \mathcal{O}_Z) \to (X, \mathcal{O}_X)$ be a morphism
of ringed spaces. Assume $i$ is a homeomorphism onto a closed subset of $X$
and $i^\sharp : \mathcal{O}_X \to i_*\mathcal{O}_Z$ is surjective.
Denote $\mathcal{I} \subset \mathcal{O}_X$ the kernel of $i^\sharp$.
The functor
$$
i_* :
\textit{Mod}(\mathcal{O}_Z)
\longrightarrow
\textit{Mod}(\mathcal{O}_X)
$$
is exact, fully faithful, with essential image those
$\mathcal{O}_X$-modules $\mathcal{G}$ such that $\mathcal{I}\mathcal{G} = 0$.
\end{lemma}

\begin{proof}
We claim that for an $\mathcal{O}_Z$-module $\mathcal{F}$ the canonical map
$$
i^*i_*\mathcal{F} \longrightarrow \mathcal{F}
$$
is an isomorphism. We check this on stalks. Say $z \in Z$ and $x = i(z)$.
We have
$$
(i^*i_*\mathcal{F})_z =
(i_*\mathcal{F})_x \otimes_{\mathcal{O}_{X, x}} \mathcal{O}_{Z, z} =
\mathcal{F}_z \otimes_{\mathcal{O}_{X, x}} \mathcal{O}_{Z, z} =
\mathcal{F}_z
$$
by Sheaves, Lemma \ref{sheaves-lemma-stalk-pullback-modules},
the fact that $\mathcal{O}_{Z, z}$ is a quotient of $\mathcal{O}_{X, x}$, and
Sheaves, Lemma \ref{sheaves-lemma-stalks-closed-pushforward}.
It follows that $i_*$ is fully faithful.

\medskip\noindent
Let $\mathcal{G}$ be a $\mathcal{O}_X$-module with
$\mathcal{I}\mathcal{G} = 0$. We will prove the canonical map
$$
\mathcal{G} \longrightarrow i_*i^*\mathcal{G}
$$
is an isomorphism. This proves that $\mathcal{G} = i_*\mathcal{F}$
with $\mathcal{F} = i^*\mathcal{G}$ which finishes the proof.
We check the displayed map induces an isomorphism on stalks.
If $x \in X$, $x \not \in i(Z)$, then $\mathcal{G}_x = 0$
because $\mathcal{I}_x = \mathcal{O}_{X, x}$ in this
case. As above $(i_*i^*\mathcal{G})_x = 0$ by
Sheaves, Lemma \ref{sheaves-lemma-stalks-closed-pushforward}.
On the other hand, if $x \in Z$, then we obtain the map
$$
\mathcal{G}_x
\longrightarrow
\mathcal{G}_x \otimes_{\mathcal{O}_{X, x}} \mathcal{O}_{Z, x}
$$
by Sheaves, Lemmas \ref{sheaves-lemma-stalk-pullback-modules} and
\ref{sheaves-lemma-stalks-closed-pushforward}. This map is an isomorphism
because $\mathcal{O}_{Z, x} = \mathcal{O}_{X, x}/\mathcal{I}_x$
and because $\mathcal{G}_x$ is annihilated by $\mathcal{I}_x$ by assumption.
\end{proof}

\begin{remark}
\label{remark-sections-support-in-closed-modules}
Let $(X, \mathcal{O}_X)$ be a ringed space. Let $Z \subset X$ be
a closed subset. For an $\mathcal{O}_X$-module $\mathcal{F}$ we can
consider the {\it submodule of sections with support in $Z$}, denoted
$\mathcal{H}_Z(\mathcal{F})$, defined by the rule
$$
\mathcal{H}_Z(\mathcal{F})(U) =
\{s \in \mathcal{F}(U) \mid \text{Supp}(s) \subset U \cap Z\}
$$
Observe that $\mathcal{H}_Z(\mathcal{F})(U)$ is a module over
$\mathcal{O}_X(U)$, i.e., $\mathcal{H}_Z(\mathcal{F})$ is an
$\mathcal{O}_X$-module. By construction $\mathcal{H}_Z(\mathcal{F})$
is the largest $\mathcal{O}_X$-submodule of $\mathcal{F}$ whose support is
contained in $Z$.
Applying Lemma \ref{lemma-i-star-equivalence} to
the morphism of ringed spaces $(Z, \mathcal{O}_X|_Z) \to (X, \mathcal{O}_X)$ we
may (and we do) view
$\mathcal{H}_Z(\mathcal{F})$ as an $\mathcal{O}_X|_Z$-module on $Z$.
Thus we obtain a functor
$$
\textit{Mod}(\mathcal{O}_X) \longrightarrow \textit{Mod}(\mathcal{O}_X|_Z),
\quad
\mathcal{F} \longmapsto \mathcal{H}_Z(\mathcal{F})
\text{ viewed as an }\mathcal{O}_X|_Z\text{-module on }Z
$$
This functor is left exact, but in general not exact.
All of the statements made above follow directly from
Lemma \ref{lemma-support-section-closed}.
Clearly the construction is compatible with the construction in
Remark \ref{remark-sections-support-in-closed}.
\end{remark}

\begin{lemma}
\label{lemma-adjoint-section-with-support}
Let $(X, \mathcal{O}_X)$ be a ringed space. Let $i : Z \to X$ be the
inclusion of a closed subset. The functor
$\mathcal{H}_Z : \textit{Mod}(\mathcal{O}_X) \to
\textit{Mod}(\mathcal{O}_X|_Z)$ of
Remark \ref{remark-sections-support-in-closed-modules}
is right adjoint to
$i_* : \textit{Mod}(\mathcal{O}_X|_Z) \to \textit{Mod}(\mathcal{O}_X)$.
\end{lemma}

\begin{proof}
We have to show that for any $\mathcal{O}_X$-module $\mathcal{F}$
and any $\mathcal{O}_X|_Z$-module $\mathcal{G}$ we have
$$
\Hom_{\mathcal{O}_X|_Z}(\mathcal{G}, \mathcal{H}_Z(\mathcal{F})) =
\Hom_{\mathcal{O}_X}(i_*\mathcal{G}, \mathcal{F})
$$
This is clear
because after all any section of $i_*\mathcal{G}$ has support in $Z$.
Details omitted.
\end{proof}










